\section{Pendahuluan}

Data time series muncul di banyak bidang, mulai dari sensor industri sampai aktivitas pengguna layanan digital. Salah satu contoh yang mudah diamati adalah jumlah pemain sebuah game pada platform Steam. Angka ini berubah setiap hari: naik saat game baru dirilis atau mendapat pembaruan besar, turun perlahan seiring waktu, dan berulang setiap pekan mengikuti kebiasaan pemain. Perubahan yang tidak biasa, misalnya lonjakan atau penurunan tajam dalam satu hari, dapat dianggap sebagai anomali yang menarik untuk dideteksi.

Rahadian dkk.~\cite{rahadian2023} menunjukkan bahwa anomali pada time series dapat dideteksi dengan mengubah potongan data menjadi gambar memakai Gramian Angular Field (GAF) dan Markov Transition Field (MTF), lalu mengklasifikasikan gambar tersebut dengan convolutional neural network (CNN). Paper tersebut juga menemukan bahwa selisih statistik matriks hasil encoding antara data normal dan anomali, terutama selisih simpangan baku (\dsd), berkorelasi sangat kuat dengan akurasi CNN, sehingga encoding yang baik dapat dipilih tanpa harus melatih CNN satu per satu.

Penelitian ini menerapkan metode yang sama pada data jumlah pemain harian enam game Steam. Tujuannya ada tiga. Pertama, memvisualisasikan data dan gambar hasil encoding untuk memperoleh insight tentang perilaku pemain. Kedua, membandingkan 14 jenis encoding dan melihat apakah temuan Rahadian dkk.\ juga berlaku pada data ini. Ketiga, memakai Grad-CAM~\cite{selvaraju2017} untuk melihat bagian gambar, dan karena itu hari, mana yang membuat CNN menganggap sebuah segmen anomali.